\begin{problem}[25 pts]
\statementfigure{%
Plot the velocity of the object, whose acceleration is shown in Figure
on the right. The velocity is zero at the moment $t=\qty{0}{\second}$.
Find the mean velocity averaged over the time interval
$[6,\qty{14}{\second}]$.
}{%
\resizebox{\linewidth}{!}{\begin{tikzpicture}
\begin{axis}[width=10cm,height=3.4cm,scale only axis,
 at={(0,0)},anchor=south west,axis lines=middle,
 xmin=0,xmax=15,ymin=-5,ymax=3.4,
 xlabel={$t$ (\unit{s})},ylabel={$a$ (\unit{\metre\per\second\squared})},
 xtick={0,2,4,6,8,10,12,14},ytick={-4,0,2},
 tick label style={font=\large},label style={font=\large},
 ylabel style={at={(axis description cs:0,1)},anchor=south west,
 rotate=0,yshift=3pt},
 grid=major,grid style={PhysicsLine},clip=false]
\addplot[physics line] coordinates
 {(0,0)(2,0)(2,2)(4,2)(4,0)(8,0)(8,2)(10,2)(10,0)
 (12,0)(12,-4)(14,-4)(14,0)};
\end{axis}
% Identical horizontal extent and plot width to velocity-blank.tikz.
\pgfresetboundingbox
\path[use as bounding box] (-.8,-.45) rectangle (11.35,4.75);
\end{tikzpicture}
}
\ifphysicsshowsolutions\else
\par\smallskip
\resizebox{\linewidth}{!}{\begin{tikzpicture}
\begin{axis}[width=10cm,height=2cm,scale only axis,
 at={(0,0)},anchor=south west,axis lines=middle,
 xmin=0,xmax=15,ymin=-0.6,ymax=3,
 xlabel={$t$ (\unit{s})},ylabel={$v$ (\unit{\metre\per\second})},
 xtick={0,2,4,6,8,10,12,14},ytick={0},
 tick label style={font=\large},label style={font=\large},
 ylabel style={at={(axis description cs:0,1)},anchor=south west,
 rotate=0,yshift=3pt},clip=false]
\addplot[only marks,mark=*,mark size=1.5pt] coordinates {(0,0)};
\end{axis}
% Identical horizontal extent and plot width to acceleration-v1.tikz.
\pgfresetboundingbox
\path[use as bounding box] (-.8,-.45) rectangle (11.35,3.35);
\end{tikzpicture}
}
\fi}
\end{problem}
\begin{solution}
\qstep{1}{Use acceleration to find the velocity changes.}
Use $v_{\rm end}=v_{\rm start}+a\,\Delta t$ on each constant-acceleration
interval. The velocity graph is horizontal where $a=0$ and has slope $a$
elsewhere.
\qstep{2}{Write the area formula before using numbers.}
For endpoint velocities $v_i,v_{i+1}$ and durations $\Delta t_i$,
add the signed trapezoid areas, then divide by the total time:
\[
\bar v=
\frac{\sum_i(v_i+v_{i+1})\,\Delta t_i}
{2\sum_i\Delta t_i}.
\]
\qstep{3}{Calculate the velocities and draw the graph.}
Starting from zero, the two positive pulses each add
$2(2)=\qty{4}{\metre\per\second}$; the last pulse adds
$-4(2)=-\qty{8}{\metre\per\second}$.
\sourcefig{0.65\linewidth}{velocity-v1}
\qstep{4}{Add the areas over 6--14 s and divide by the time.}
The velocity at 6 s is also $\qty{4}{\metre\per\second}$.
The four areas give
\[
\Delta x=(8+12+16+8)\,\unit{\metre}=\qty{44}{\metre},\qquad
\bar v=\frac{\qty{44}{\metre}}{\qty{8}{\second}}
=\qty{5.50}{\metre\per\second}.
\]
\end{solution}

\quiznextproblem
\begin{problem}[25 pts]
The tram starts moving with a constant acceleration of
\qty{0.4}{\metre\per\second\squared}. How long can a passenger stay after
the beginning of the tram movement, then catch up to him, running at a
constant speed of \qty{10}{\metre\per\second}?
\end{problem}
\begin{solution}
\qstep{1}{Choose the model and sketch the limiting meeting.}
Assume a shared starting position, a tram starting from rest, and a
passenger who waits there for time $t_{\rm w}$ before instantly reaching running
speed $v_{\rm P}$. Interpret the question as the longest possible wait.
\sourcefig{0.48\linewidth}{tram-limit}
At the limit, equal positions and speeds make the two curves tangent.
\qstep{2}{Write the position equations.}
With $+x$ along the motion and $t\geq t_{\rm w}$,
\[
x_T=\frac{at^2}{2},\qquad x_P=v_{\rm P}(t-t_{\rm w}).
\]
Before $t_{\rm w}$, the passenger stays at $x=0$.
\qstep{3}{Require a physically possible meeting.}
Equating positions gives $at^2-2v_{\rm P}t+2v_{\rm P}t_{\rm w}=0$, hence
\[
t_\pm=\frac{v_{\rm P}\pm\sqrt{v_{\rm P}^2-2av_{\rm P}t_{\rm w}}}{a}.
\]
For $a,v_{\rm P}>0$ and a positive wait, a meeting after the wait exists when
$v_{\rm P}^2-2av_{\rm P}t_{\rm w}\geq0$. Therefore
\[
t_{\rm w,max}=\frac{v_{\rm P}}{2a},\qquad 0\leq t_{\rm w}\leq\frac{v_{\rm P}}{2a}.
\]
\qstep{4}{Calculate and interpret the answer.}
\[
t_{\rm w,max}=\frac{\qty{10}{\metre\per\second}}
{2(\qty{0.4}{\metre\per\second\squared})}=\qty{12.5}{\second}.
\]
At the limit $t_*=\frac{v_{\rm P}}{a}=\qty{25.0}{\second}$.
For a smaller positive delay, $t_-$ is the first meeting; $t_+$ is a
later tram recatch if both continue. Zero delay means initial coincidence,
not a chase after waiting.

\textit{Alternative:} minimize the separation
$d(t)=\frac{at^2}{2}-v_{\rm P}(t-t_{\rm w})$.
Its minimum at $t=\frac{v_{\rm P}}{a}$ is $v_{\rm P}t_{\rm w}-\frac{v_{\rm P}^2}{2a}$;
requiring this to be nonpositive gives the same bound.
\end{solution}

\quizpairbreak
\begin{problem}[25 pts]
\statementfigure{%
A drop of rain falls at an angle of $\alpha_1=\ang{30}$ to the vertical
axis when the wind speed is $u_1=\qty{10}{\metre\per\second}$.
What is the angle $\alpha_2$ at which the drop falls after the wind has
been increased to $u_2=\qty{20}{\metre\per\second}$?
Assume that the velocity of the raindrops with respect to the air is constant.
}{\resizebox{0.65\linewidth}{!}{\begin{tikzpicture}[font=\small]
\draw[physics line] (0,0)--(4,0);
\draw[dashed,-{Latex}] (2.9,2.8)--(2.9,0.25)
 node[pos=.4,right]{Vertical};
\draw[physics line,-{Latex}] (2.9,2.8)--(1.6,0.2);
\draw (2.9,2.1) arc[start angle=270,end angle=243.4,radius=.7];
\node at (2.63,1.85) {$\alpha$};
\draw[physics accent,-{Latex}] (2.8,3.2)--(1.2,3.2)
 node[midway,above]{Wind};
\node at (2,-.4) {Schematic};
\end{tikzpicture}
}}
\end{problem}
\begin{solution}
\qstep{1}{Draw the velocities and state the working assumption.}
Assume vertical downward motion relative to the air and horizontal wind
in the same direction in both states. Measure the observed angle in the
ground frame. Take $+y$ downward; $v_y$ is the unchanged vertical velocity component.
\sourcefig{0.49\linewidth}{rain-vectors}
The horizontal leg changes; the vertical leg does not.
\qstep{2}{Relate each angle to the components.}
Velocity addition gives horizontal component $u_i$ and vertical component $v_y$.
Because the angle is measured from the vertical,
\[
\tan\alpha_1=\frac{u_1}{v_y},\qquad
\tan\alpha_2=\frac{u_2}{v_y}.
\]
\qstep{3}{Eliminate the unknown vertical speed.}
Using $v_y=\frac{u_1}{\tan\alpha_1}$ gives
\[
\alpha_2=\arctan\!\left(\frac{u_2}{u_1}\tan\alpha_1\right).
\]
The downward-drift solution is the acute angle.
\qstep{4}{Evaluate the angle.}
\[
\alpha_2=\arctan\!\left(2\tan30^\circ\right)
=\arctan\!\left(\frac{2}{\sqrt3}\right)\approx\ang{49.1}.
\]
The same relation follows from two right velocity triangles with a
common vertical leg. Without the vertical air-relative-motion assumption,
the source wording does not fix a unique numerical angle.
\end{solution}

\quiznextproblem
\begin{problem}[25 pts]
A particle moves in uniform circular motion over a horizontal $xy$ plane.
It takes \qty{20}{\second} to make a full circle. At one instant, it moves
through the point with coordinates $[\qty{4}{\metre},\qty{4}{\metre}]$
with the acceleration of
$(0\hat{\mathbf i}+8\hat{\mathbf j})\,\unit{\metre\per\second\squared}$.
Find the $x$ and $y$ coordinates of the center of the circular path.
\end{problem}
\begin{solution}
\qstep{1}{Use the direction of centripetal acceleration.}
\sourcefig{0.37\linewidth}{circle-v1}
Centripetal acceleration points toward the center.
Here it is vertically upward, so $x_C=x$ and $y_C=y+R$.
\qstep{2}{Find the radius from centripetal acceleration.}
Let $a_c=|\mathbf a|$ be its magnitude and $v$ the constant speed.
One revolution covers $2\pi R$ in time $T$, so
\[
a_c=\frac{v^2}{R},\qquad v=\frac{2\pi R}{T}
\quad\Longrightarrow\quad a_c=\frac{4\pi^2R}{T^2}.
\]
\qstep{3}{Write the symbolic answer.}
\[
R=\frac{a_cT^2}{4\pi^2},\qquad
x_C=x,\qquad y_C=y+\frac{a_cT^2}{4\pi^2}.
\]
\qstep{4}{Calculate the radius and coordinates.}
\[
R=\frac{8(20)^2}{4\pi^2}\,\unit{\metre}
=\frac{800}{\pi^2}\,\unit{\metre}.
\]
\[
x_C=\qty{4}{\metre},\qquad
y_C=\left(4+\frac{800}{\pi^2}\right)\unit{\metre}.
\]
\textbf{Answer:} $(x_C,y_C)\approx(4,85.1)\,\unit{\metre}$.
The radius is added to the point's height; it is not the center's height itself.

\textit{Vector alternative:} with $\omega=\frac{2\pi}{T}$,
$\mathbf a=\omega^2(\mathbf r_C-\mathbf r)$ gives
$\mathbf r_C=\mathbf r+\frac{T^2}{4\pi^2}\mathbf a$, the same result.
\end{solution}
